\section{Introduction}

The introduction argues that on-device models need a way to learn on the job, that the hard part is deciding what to learn from, and that we measure this over long lives rather than single updates.

\begin{itemize}
\tightlist
\item
  \textbf{Motivation:} small models on phones and devices meet new tasks after release; retraining in the cloud is costly and slow. A model that improves from its own daily work would close that gap.
\item
  \textbf{The core problem:} learning on the job needs a judge of which episodes were good. Without a true outcome, the options are the model's own judgment or the user's reaction, and both can be wrong.
\item
  \textbf{The gap:} most work measures one round of improvement on a fixed benchmark. Little measures what happens over weeks of nightly updates, where small errors in the judge compound.
\item
  \textbf{Our approach:} a ``creature'' that works by day and sleeps on its lessons at night, run for 14 days to a year under different judges, including a separate frozen model, the feedback sense, that reads the user's reaction. We call what it must provide discernment: telling an informed, clear verdict on the result from a polite, vague or uninformed reaction. The other judges are simplified versions of published approaches, compared with it under the same model, tasks and training budget.
\item
  \textbf{Contributions:}

  \begin{enumerate}
  \def\labelenumi{\arabic{enumi}.}
  \tightlist
  \item
    A long-horizon test of nightly self-training for small models, on WebShop and $\tau^2$-bench, with an oracle as the upper bound.
  \item
    Evidence that learning from the model's own judgment decays, and that what decides a life is the share of failed attempts among its lessons, not their number ($\rho$ = -0.81 across 97 lives): a few accepted give-ups can collapse a life.
  \item
    A feedback sense with discernment, which first asks whether the user gave a clear verdict. It removes most bad lessons and reaches near-oracle scores for two weeks, then fades for reasons we leave open.
  \item
    Evidence that learning from reactions fails when users cannot see the outcome, and that even correct but lopsided lessons cause forgetting.
  \end{enumerate}
\end{itemize}
